\documentclass[11pt]{article}
\usepackage[T1]{fontenc}
\usepackage{lmodern,amsmath,amssymb,amsthm,geometry,enumitem,graphicx,xcolor,tikz,hyperref,fancyhdr}
\geometry{margin=0.9in}
\definecolor{courseblue}{HTML}{285E80}
\hypersetup{colorlinks=true,linkcolor=courseblue,urlcolor=courseblue}
\pagestyle{fancy}\fancyhf{}\fancyhead[L]{MATH\& 142 -- Precalculus II}\fancyhead[R]{Dr. Lee}\fancyfoot[C]{\thepage}
\newtheoremstyle{course}{7pt}{7pt}{}{}{\bfseries\color{courseblue}}{.}{.5em}{}
\theoremstyle{course}\newtheorem{definition}{Definition}[section]\newtheorem{example}{Example}[section]
\newenvironment{practice}{\par\medskip\noindent\textbf{Try it. }}{\par\medskip}
\setlist[enumerate]{itemsep=5pt}

\newcommand{\workspace}{\par\vspace{.9in}}
\begin{document}
\begin{center}{\LARGE MATH\& 142 --- Comprehensive Final Examination}\\[4pt]
{\large Original assessment, Modules 1--10}\end{center}
Name:\ \rule{3in}{.4pt}\hfill Date:\ \rule{1.35in}{.4pt}

\smallskip\noindent\textbf{Instructions:} Show work and justify results. Supply exact answers unless a decimal approximation is requested. Use radians unless a degree symbol appears. Ten questions, 10 points each, total 100 points.
\begin{enumerate}
\item \textbf{Radian measure and the unit circle.} Convert $240^\circ$ to radians. Find $\sin(4\pi/3)$, $\cos(4\pi/3)$, and $\tan(4\pi/3)$ exactly.\workspace
\item \textbf{Graph transformations.} For $y=-2\sin(3(x-\pi/6))+1$, determine amplitude, period, phase shift, midline, and range.\workspace
\item \textbf{Inverse functions and identities.} Evaluate $\arcsin(\sin(5\pi/6))$. Prove $(1-\sin^2x)/\cos x=\cos x$ on its domain.\workspace
\item \textbf{Trigonometric equations.} Solve $\cos(2x)=1/2$ for all $x$ in $[0,2\pi)$.\workspace
\item \textbf{An oblique triangle.} Two sides of a triangle have lengths $6$ and $10$, with included angle $120^\circ$. Calculate the length of the opposite side exactly.\workspace
\item \textbf{Vectors and dot products.} For $\mathbf u=\langle3,4\rangle$ and $\mathbf v=\langle4,-3\rangle$, determine $\mathbf u+\mathbf v$, both magnitudes, and the angle between the vectors.\workspace
\item \textbf{Complex numbers and rotation.} Let $z=2(\cos30^\circ+i\sin30^\circ)$ and $w=3(\cos60^\circ+i\sin60^\circ)$. Find the product in polar and rectangular forms.\workspace
\item \textbf{Parametric curves and conics.} Eliminate $t$ from $x=4\cos t$, $y=3\sin t$. Identify the resulting conic, its semiaxes, and its foci.\workspace
\item \textbf{Sequences and series.} A geometric sequence has $a_1=9$ and ratio $q=1/3$. Find $a_6$ and the infinite sum $\sum_{n=1}^{\infty}a_n$.\workspace
\item \textbf{Course synthesis.} A point starts at $z_0=4$ and obeys $z_{n+1}=\frac12 i z_n$. Determine $z_1,z_2,z_3$, then give a formula for $z_n$. Explain the geometric pattern, including the behavior of magnitude and direction.\workspace
\end{enumerate}
\end{document}
